\begin{claim}{D.2.57}{Proved}{The bottom sigma group survives through the fourteenth stage}
The two remaining incoming differentials in D.2.55 vanish:
\[
 d_6:E^6_{13,2}\longrightarrow E^6_{7,7},\qquad
 d_7:E^7_{14,1}\longrightarrow E^7_{7,7},\qquad
                         E^8_{7,7}=\ZZ/4.
 \tag{D.2.57a}
\]
In an actual fourteen-skeleton $P_{14}$ of
$L=Y\langle7\rangle$, the image of the bottom stable
element $2\sigma$ is nonzero.  It is detected in Adams
filtration two by the image of $h_0h_3$.

Consequently the detecting class of D.2.54 has the exact
remaining obstruction
\[
       d_8\omega=2n\bar\alpha\in\ZZ/4,
       \qquad
       h_Ml_4=0\quad\Longleftrightarrow\quad
                   2n\bar\alpha\in d_8(H_0).
       \tag{D.2.57b}
\]
The integer $n$ and the source-changing image $d_8(H_0)$
are not evaluated here.  In particular this finite-stage
survival result does not decide the actual full-target
obstruction or D.2.4.
\end{claim}
\begin{proof}
\textbf{A four-dimensional Steenrod submodule.}
Let $a_7\in H^7(L;\FF_2)$ be the nonzero bottom class, and put
\[
 s_9=\operatorname{Sq}^2a_7,\qquad
 n_{11}=\operatorname{Sq}^4a_7,\qquad
 z_{13}=\operatorname{Sq}^4\operatorname{Sq}^2a_7.
 \tag{D.2.57h1}
\]
Through degree fourteen these four nonzero classes span the entire
Steenrod cyclic submodule generated by $a_7$. Its only nonzero positive
operations in this range are the three in \textup{(D.2.57h1)}.

Here is a direct verification in the actual spaces. By D.2.56,
double cohomology suspension embeds $H^{\le14}L$ into $H^{\le12}E$.
Under this embedding, the four classes in
\textup{(D.2.57h1)} are the coefficient functionals of the outer
restricted generators
\[
 \xi a,\qquad \xi b,\qquad \xi c,\qquad \xi d
 \quad\hbox{in degrees }5,7,9,11\hbox{ of }H_*T,
 \tag{D.2.57h2}
\]
restricted to $H_*E$. Here $a,b,c,d$ are the homology labels dual
to $u_B,v_B,u_B^2,w_B$.
The top Nishida formulas
\cite[Chapter III, Theorem 1.3(3), pp.~217--218]{I_CohenLadaMay1976}
give these identifications. In particular the only linear-label
inputs relevant to the indicated even operations are
$\operatorname{Sq}_*^1b=a$,
$\operatorname{Sq}_*^2c=a$, and
$\operatorname{Sq}_*^2d=b$.
The other degree-five inputs, the label dual to $u_Bv_B$ and
$\xi a$, have zero second square.

Every odd homology square has zero coefficient on an outer restricted
generator. Indeed, for a restricted generator the odd top Nishida
formula has only Pontryagin-product and Browder-bracket terms.
The restricted adjoint identity expresses brackets of restricted
words as ordinary Lie words, so these have no outer restricted-power
component. On labels, ordinary Hall words, and Pontryagin products,
the same assertion follows from naturality and the two Cartan rules.
Consequently
\[
 \operatorname{Sq}^1a_7=\operatorname{Sq}^1s_9
 =\operatorname{Sq}^1n_{11}
 =\operatorname{Sq}^1z_{13}=0.
 \tag{D.2.57h3}
\]
Adem gives $\operatorname{Sq}^2s_9=0$ and
\[
 \operatorname{Sq}^2n_{11}
 =\operatorname{Sq}^6a_7+\operatorname{Sq}^5
       \operatorname{Sq}^1a_7=0.
 \tag{D.2.57h4}
\]
For the last equality, $\operatorname{Sq}^6a_7=0$ follows after
double suspension from instability on the degree-five class; the
injectivity in degree thirteen is part of D.2.56.
Adem, \textup{(D.2.57h3)--(D.2.57h4)}, and instability now give
all the remaining vanishings through degree fourteen.

The established low-degree tables show that $a_7,s_9,n_{11}$ are
nonzero. To check $z_{13}$ without a full high-degree table, let
$P_{4v}$ be the coefficient of $u_B^4v_B$ in
$\mathcal G^{-1}v_B\partial\mathcal G/\partial v_B$, where
$\mathcal G=\sum_m[m]m$ is the group-like label series. This is a
primitive homology class. Set
\[
 w_{11}=\xi d+Z_{11}+P_{4v},\qquad
 Z_{11}=[u_Bt_B]+a[t_B].
 \tag{D.2.57h5}
\]
All three summands are primitive classes of the fiber $F$. The only
polynomial generators of $H^*K(\FF_2,4)$ in degree eleven are
$x_{11}=\operatorname{Sq}^5\operatorname{Sq}^2x_4$ and
$y_{11}=\operatorname{Sq}^4\operatorname{Sq}^2
\operatorname{Sq}^1x_4$. Their pairing vectors on the three
summands in \textup{(D.2.57h5)} are
\[
 (1,0),\qquad(0,1),\qquad(1,1),
 \tag{D.2.57h6}
\]
respectively. For the last pair, Cartan on the finite label series
gives $\beta P_{4v}=P_p$ and
$\operatorname{Sq}_*^4P_{4v}=P$, with the previously specified
degree-ten Newton class $P_p$ and degree-seven class $P$.
The earlier tables give the first two pairs. Thus $w_{11}$
annihilates both degree-eleven base generators. Primitivity also
annihilates every positive cohomology product, so the exact edge
description in D.2.56 places $w_{11}$ in $H_{11}E$. Finally,
\[
 \operatorname{Sq}_*^4w_{11}=\xi b+P+Q=S,
 \qquad\operatorname{Sq}_*^2S=\xi a.
 \tag{D.2.57h7}
\]
Hence $z_{13}$ evaluates to one on the actual double suspension
of $w_{11}$, proving its nonvanishing.

Restricting to the ordinary fourteen-skeleton $P_{14}$ preserves this
submodule: cohomology restriction is an isomorphism through degree
thirteen and injective in degree fourteen. No minimal CW model is
required. Dually, $H_*P_{14}$ has a four-dimensional quotient
comodule, with basis denoted $a,s,n,z$ in degrees $7,9,11,13$ and
reduced left Milnor coaction
\[
 \begin{aligned}
 \overline\psi(a)&=0,&
 \overline\psi(s)&=\xi_1^2\otimes a,\\
 \overline\psi(n)&=\xi_1^4\otimes a,&
 \overline\psi(z)&=\xi_1^4\otimes s+\xi_2^2\otimes a.
 \end{aligned}
 \tag{D.2.57h8}
\]
The conventions are
$\Delta\xi_j=\sum_i\xi_{j-i}^{2^i}\otimes\xi_i$ and the ordinary
left coaction dual to Steenrod operations.


\medskip\noindent
\textbf{Restriction to an actual finite skeleton.}
Choose the ordinary finite fourteen-skeleton $P_{14}$
as above, under its specified bottom sphere.  The restriction
\[
 H^i(L;\FF_2)\longrightarrow H^i(P_{14};\FF_2)
 \tag{D.2.57c}
\]
is an isomorphism for $i\le13$ and is injective for
$i=14$.  As explained in \textup{(D.2.57h8)}, this
gives the required actual quotient without identifying
all of $H^{14}(P_{14})$ with $H^{14}(L)$.

Let $\mathcal A$ be the mod-two Steenrod algebra.
Dualizing the four-class submodule gives a left
$\mathcal A_*$-comodule quotient of $H_*P_{14}$,
with basis $A,S,N,Z$ in degrees $7,9,11,13$.
Its reduced coaction is
\[
 \overline\psi A=0,\qquad
 \overline\psi S=\xi_1^2\otimes A,\qquad
 \overline\psi N=\xi_1^4\otimes A,\qquad
 \overline\psi Z=\xi_1^4\otimes S+\xi_2^2\otimes A.
 \tag{D.2.57d}
\]
Here $\xi_i$ are the standard Milnor generators, with
\[
 \Delta\xi_k=\sum_{i=0}^k\xi_{k-i}^{2^i}\otimes\xi_i,
                 \qquad\xi_0=1;
 \tag{D.2.57e}
\]
see \cite[Theorem 3, p.~162, and Appendix 1, p.~171]
{EH5Milnor1958Steenrod}.
For example $\xi_2^2$ in \textup{(D.2.57d)} is
forced by the $\operatorname{Sq}^4\operatorname{Sq}^2$
value and coassociativity; the possible $\xi_1^6$
coefficient is zero because $\operatorname{Sq}^6a_7=0$.

\medskip\noindent
\textbf{A complete finite cobar certificate.}
Use the normalized cobar complex
$\overline{\mathcal A}_*^{\otimes s}\otimes
\langle A,S,N,Z\rangle$.  Its differential splits
each algebra tensor by the reduced coproduct and
appends a reduced comodule term on the right.
This convention corresponds to the left coaction in
\textup{(D.2.57d)}.  Write $[u|v]M$ for a basis
element of $C^{2,16}$.  Then
\[
                         t=[\xi_1|\xi_1^8]A
                         \tag{D.2.57f}
\]
is a cycle: both algebra factors are primitive and
$A$ has zero reduced coaction.

Define a functional $\Phi:C^{2,16}\to\FF_2$ to be
one precisely on the following eleven basis elements,
and zero on every other basis element:
\[
\begin{array}{c|cc}
 M&u&v\\ \hline
 A&\xi_1&\xi_1^8\\
 A&\xi_3&\xi_1^2\\
 A&\xi_1^4\xi_2&\xi_1^2\\
 A&\xi_1^7&\xi_1^2\\
 A&\xi_1^2\xi_2^2&\xi_1\\
 A&\xi_1^8&\xi_1\\
 S&\xi_1^3&\xi_1^4\\
 S&\xi_1^5&\xi_1^2\\
 S&\xi_2^2&\xi_1\\
 S&\xi_1^6&\xi_1\\
 Z&\xi_1&\xi_1^2
\end{array}
 \tag{D.2.57g}
\]
There are exactly thirteen incoming basis elements in
$C^{1,16}$.  The next table lists all of them, together
with the number of terms of their differential on which
$\Phi$ is one.  Every number is even.
\[
\begin{array}{c|c|r}
 M&u\text{ in }[u]M&\text{number}\\ \hline
 A&\xi_2^3&2\\
 A&\xi_1^2\xi_3&2\\
 A&\xi_1^3\xi_2^2&2\\
 A&\xi_1^6\xi_2&2\\
 A&\xi_1^9&2\\
 S&\xi_3&2\\
 S&\xi_1\xi_2^2&2\\
 S&\xi_1^4\xi_2&2\\
 S&\xi_1^7&4\\
 N&\xi_1^2\xi_2&0\\
 N&\xi_1^5&0\\
 Z&\xi_2&0\\
 Z&\xi_1^3&2
\end{array}
 \tag{D.2.57h}
\]
This is a direct expansion of
\textup{(D.2.57d)--(D.2.57e)}.  For instance the two
supported terms in $d[\xi_3]S$ are
$[\xi_2^2|\xi_1]S$ and $[\xi_3|\xi_1^2]A$;
those in $d[\xi_1^3]Z$ are
$[\xi_1|\xi_1^2]Z$ and $[\xi_1^3|\xi_1^4]S$.
Thus $\Phi$ annihilates the entire incoming boundary
space, whereas $\Phi(t)=1$.  In particular
\[
       [t]\ne0\quad\hbox{in the quotient-comodule
                     cobar group }\operatorname{Ext}^{2,16}.
       \tag{D.2.57i}
\]
The comodule quotient preserves the bottom class.
Therefore the image of $h_0h_3$ in
$\operatorname{Ext}_{\mathcal A}^{2,16}
(H^*P_{14},\FF_2)$ is nonzero as well.

The authored exact calculation
\texttt{code/EH5\_bottom\_sigma\_cobar.py} checks
the same thirteen columns, the left coaction identity,
and $d^2=0$; its four-class input and explicit dual
certificate are saved in the corresponding
\texttt{code/outputs/EH5\_bottom\_sigma\_four\_class\_*}
files.  The tables above provide the complete
nonboundary proof independently of that implementation.
The larger operation-matrix calculation is a separate
check and is unnecessary for this implication.

\medskip\noindent
\textbf{Adams permanence and the two incoming differentials.}
The Hopf-invariant-one element $\sigma$ is detected
by $h_3$: its cone has nonzero eighth square between
the two cells.  Multiplication by two is detected by
$h_0$, and the nonzero product $h_0h_3$ detects
$2\sigma$.  The survival of this product is also
certified by its nonzero image under the sphere-to-$j$
unit in \cite[Theorem 1.2 and Remark 1.3, p.~3]
{EH5BrunerRognes2022}.
The bottom sphere map sends its cobar representative to
\textup{(D.2.57f)}.  We use the natural Adams spectral
sequence and its comodule-resolution identification in
\cite[Part III, Theorem 15.1, pp.~370--371, and
pp.~372--377]{EH5Adams1974AHSS}.
Naturality from the permanent sphere class prevents
the nonzero image in \textup{(D.2.57i)} from supporting
an Adams differential.  It cannot be hit either.
The only possible incoming differential to bidegree
$(2,16)$ is $d_2$ from $(0,15)$, whose group is zero
because $H_{15}(P_{14};\FF_2)=0$.
Higher incoming differentials have negative source
filtration.

Adams convergence for this finite connective spectrum
therefore detects a nonzero completed stable image
of the specified $2\sigma$.  Its two-local stable
image is consequently nonzero as well: a zero element
would have zero completed image.  Thus
\[
               2i_7\sigma\ne0\quad\hbox{in }\pi_{14}^sP_{14}.
               \tag{D.2.57j}
\]
By D.2.55 the bottom subgroup through dimension twelve
is exactly $\ZZ/4$.  Every nonzero subgroup of
$\ZZ/4$ contains twice its generator.  Consequently
\textup{(D.2.57j)} forces this entire bottom subgroup
to inject through dimensions thirteen and fourteen.
In the cellular exact couple this says precisely that
the incoming $d_6$ and $d_7$ at $(7,7)$ are zero.
Together with D.2.54--D.2.55 it proves
\textup{(D.2.57a)}.  Substituting the now exact target
in D.2.54 gives \textup{(D.2.57b)}, while retaining
all its source-changing ambiguity.
\end{proof}
